\chapter{总结与展望}
\label{chap:conclusion}

\section{全文总结}

大规模低轨（LEO）卫星巨型星座凭借其低时延、广覆盖的通信能力，正在成为下一代天地一体化信息网络的战略基础设施。然而，星间链路（ISL）带宽有限、卫星轨道参数完全公开的特性，使 LEO 星座网络面临链路洪泛攻击（LFA）的严峻威胁。本文围绕"如何在星载可编程数据面的严苛 SWaP 约束下，构建低时延、低资源消耗且能有效对抗 LFA 退化策略的在网安全防御体系"这一核心问题，以在网计算为技术路线，开展了从单节点多签名检测到多节点协同推回缓解的系统性研究。全文的主要研究成果总结如下。

（1）\textbf{面向 LEO 网络的多签名在网检测与队列缓解方法（MS-SatShield）。} 本文在第三章中针对单节点在 P4 数据面上进行 LFA 检测时面临的特征退化困境，提出了一种多签名在网检测与队列缓解方法。该方法的核心创新在于设计了级联式两级检测流水线架构：第一级采用“Count-Min Sketch 频率估计 + 固定容量候选缓存”的 sketch-based 候选生成机制，以固定内存维护全局近似频次并显式恢复候选 key 集合；第二级在候选集上部署双缓冲位图结构，在线提取扇入度（FI）、扇出度（FO）和历史持续性（Persist）等多维结构特征签名。通过“仅对候选集合维护 FI/FO”的设计，该方法将 distinct 统计的状态成本严格控制在 $O(K \cdot m)$ 范围内（默认 $K=1000$, $m=256$, 数据面总状态约 200~KB），天然契合星载平台的 SWaP 约束。实验结果验证了该方法的有效性：在多 decoy 分摊退化设定下（$M=100$），仅靠聚合速率的基线方案 F1 降至 0.017，而引入 FI/FO 后的 MS-SatShield F1 恢复至 0.908，FI/FO 特征提供了核心检测增益；在脉冲式 on-off 攻击下，持续性特征通过加性奖励与候选集扩展机制将 Recall 提升至 0.999，良性吞吐恢复至 0.885；消融实验表明了 $K \approx 1000$、$m = 256$、$\Delta = 1.0$~s 在检测精度与资源开销之间经验上较优的参数折中。

（2）\textbf{基于抑制型扩散的多星协同在网计算与推回缓解方法（CoopSatShield）。} 本文在第四章中针对单星检测视野的局限性与"被动就地丢弃"模式的不足，提出了一种多星协同在网计算与逐跳推回缓解方法。该方法包含两个核心组件：抑制型流言传播协议（Suppressed Gossip Protocol）和逐跳推回缓解机制（Hop-by-hop Pushback）。前者通过事件驱动、区域限制（ROI, $R=4$ 跳）与基于 Trickle 思想的冗余抑制三重约束，在严格限制信令开销的前提下实现局部区域内多颗卫星之间的攻击证据快速交叉验证与聚合；后者利用入端口贡献分布实现无需全局逆路由的推回令牌传播，将攻击抑制点从受害链路前推至上游入口，有效减少攻击流在网络中的无效传输。实验结果表明了该方法的显著优势：在 $M=100$ 退化下，CoopSatShield 将路径积分意义下的无效带宽占用成本（$W_{\text{waste}}$）降至无防御方案的 11.0\%，相比 MS-SatShield 进一步降低 6.1 个百分点；攻击期间良性吞吐稳态维持在 $0.959$，相比 MS-SatShield 提升 7.3 个百分点；推荐配置 $R=4, k=3$ 下，按保守 $n=100$ 口径统计的控制带宽仅占 ISL 容量的 0.004\%，抑制机制使 EG 消息数比无抑制基线减少 45\%，证实了协同通信不会引发信令风暴。

上述两项研究成果共同构成了面向大规模 LEO 卫星网络的在网计算安全防御体系，实现了从"单节点在网计算"到"多节点协同在网计算"的完整技术链条。在理论层面，本文揭示了 LFA 攻击退化策略与多维特征签名之间的对偶约束关系——攻击者无法在压缩速率特征的同时压缩 FI/FO 特征，为多签名检测提供了理论基础。在方法层面，本文提出了级联式两级检测架构和抑制型邻居扩散协同框架，将安全检测与缓解逻辑完整下沉至星载可编程数据面，规避了星地控制环路的时延瓶颈。在系统层面，数据面总状态开销控制在约 200~KB，控制信令带宽不超过 ISL 容量的 0.004\%，满足了星载环境的资源约束。

\section{未来展望}

本文围绕低轨卫星网络中的链路洪泛攻击防御开展了研究，并在单星检测、多星协同和逐跳推回等方面取得了较好的实验结果。随着星座能力持续提升，未来的研究还需要考虑新型业务载荷和更复杂网络组织形态带来的新需求。基于这一判断，后续工作可重点从以下两个方向继续展开。

（1）\textbf{面向多业务与在轨计算载荷的防御扩展。} 未来低轨卫星网络将不再只是单纯的宽带转发网络，而会同时承载公众接入、企业专网、政府业务以及在轨 AI 处理和空间侧算力协同等任务\citess{starlink_direct_to_cell_2026, amazon_leo_enterprise_2025, xinhua_space_computing_2026}。这意味着网络中的关键对象将从传统链路和队列进一步扩展到模型分发路径、在轨推理结果、跨星算力调度和业务切片保障等新型载荷。相应地，后续研究可在本文框架基础上进一步考虑面向多类业务与计算任务的分级检测、差异化缓解与联合保障机制，使在网防御能够同时适应通信业务和计算载荷并存的运行环境。

（2）\textbf{面向跨区域与异构星座场景的协同扩展。} 本文第四章主要围绕单一星座内的局部协同与逐跳推回展开，而未来低轨网络很可能朝着更大规模、跨区域互联、多轨道层协作以及异构星载平台并存的方向发展。在这类场景中，攻击证据和缓解动作不一定局限于单一 ROI 或单一星座内部，而可能涉及跨轨道面、跨区域甚至跨星座的联合响应。因此，后续可进一步研究更大范围协同条件下的证据传播、区域间推回衔接和异构数据面协同机制，将本文提出的多签名检测与推回缓解扩展到更复杂的网络组织形态中，以提升其在未来空间信息网络中的适用性和推广价值。
